% ---------------------------------------------------------------------------
% Figure preamble for the short lecture notes.
%
% Loaded by main.tex ONLY, deliberately NOT by preamble.tex, so that anything
% reusing preamble.tex does not inherit the whole figure machinery as a side
% effect. Reuse is opt-in: \input this file explicitly.
%
% Every schematic in src/figures/ is drawn against the styles defined here.
% Do not hard-code colours or line widths in a figure file; add a style here
% instead, so the whole set stays visually consistent.
% ---------------------------------------------------------------------------

\usepackage{tikz}
\usepackage{pgfplots}
% adjustbox: shrinks a figure to the text width only when it would otherwise
% overflow, so a schematic that runs a little wide is scaled instead of spilling
% into the margin. Figures narrower than \linewidth are left completely alone.
\usepackage{adjustbox}
\pgfplotsset{compat=1.16}
\usetikzlibrary{
  arrows.meta,
  positioning,
  calc,
  decorations.pathreplacing,
  decorations.markings,
  shapes.geometric,
  fit,
  backgrounds,
  patterns
}

% --- Palette -----------------------------------------------------------------
% One accent, one "broken/anomalous" colour, one muted grey. Every colour
% distinction is ALSO carried by a shape, weight or dash difference, so the
% figures survive greyscale printing.
\definecolor{aqmaccent}{RGB}{31,78,121}   % ~blue!45!black, matches notebookbox frame
\definecolor{aqmbroken}{RGB}{168,42,34}   % ~red!70!black, reserved for broken/asymmetric/anomalous
\definecolor{aqmmuted}{RGB}{125,125,125}  % de-emphasis
\definecolor{aqmfill}{RGB}{222,233,243}   % pale accent wash for filled regions
\definecolor{aqmfillb}{RGB}{246,225,222}  % pale broken-colour wash

% --- Shared metrics ----------------------------------------------------------
\newcommand{\aqmlinewidth}{0.9pt}
\newcommand{\aqmthinwidth}{0.5pt}

% --- Node / flowchart styles -------------------------------------------------
\tikzset{
  aqm/.style={line width=\aqmlinewidth, line cap=round, line join=round},
  aqmbox/.style={
    aqm, draw=black, fill=white, rounded corners=1.5pt,
    inner sep=4pt, align=center, font=\small
  },
  aqmboxaccent/.style={aqmbox, draw=aqmaccent, fill=aqmfill},
  aqmboxbroken/.style={aqmbox, draw=aqmbroken, fill=aqmfillb},
  aqmdiamond/.style={
    aqm, draw=black, fill=white, diamond, aspect=1.7,
    inner sep=2pt, align=center, font=\small
  },
  aqmarrow/.style={aqm, draw=black, -{Stealth[length=5pt,width=4pt]}},
  aqmarrowaccent/.style={aqmarrow, draw=aqmaccent},
  aqmarrowbroken/.style={aqmarrow, draw=aqmbroken, dashed},
  aqmarrowmuted/.style={aqmarrow, draw=aqmmuted, line width=\aqmthinwidth},
  % Lattices
  aqmsite/.style={
    aqm, circle, draw=black, fill=aqmaccent, minimum size=5pt, inner sep=0pt
  },
  aqmsiteopen/.style={aqmsite, fill=white},
  aqmsitealt/.style={aqmsite, fill=white, draw=black, line width=\aqmlinewidth},
  aqmbond/.style={aqm, draw=black},
  aqmbondweak/.style={aqm, draw=black, line width=\aqmthinwidth, dash pattern=on 2pt off 1.5pt},
  aqmbondstrong/.style={aqm, draw=black, line width=1.8pt},
  aqmcell/.style={aqm, draw=aqmmuted, line width=\aqmthinwidth, dashed, rounded corners=1pt},
  % Text
  aqmlab/.style={font=\small, inner sep=1.5pt},
  aqmlabsmall/.style={font=\footnotesize, inner sep=1pt},
  aqmlabmuted/.style={aqmlabsmall, text=aqmmuted},
  aqmpanel/.style={font=\small\bfseries, inner sep=2pt},
  aqmbrace/.style={aqm, decorate, decoration={brace, amplitude=4pt}, draw=aqmmuted,
                   line width=\aqmthinwidth},
  % TikZ-level twin of the pgfplots "aqmguide" style below, so the same dashed
  % construction-line look is available inside a plain tikzpicture.
  aqmguide/.style={draw=aqmmuted, line width=0.6pt, dash pattern=on 2pt off 1.5pt}
}

% --- pgfplots styles ---------------------------------------------------------
% aqmplot: the default look for every dispersion / DOS panel in the set.
\pgfplotsset{
  aqmplot/.style={
    width=6cm, height=4.2cm,
    axis lines=middle,
    axis line style={aqm, black, -{Stealth[length=4pt,width=3pt]}},
    tick align=outside,
    tick style={black, line width=\aqmthinwidth},
    label style={font=\small},
    tick label style={font=\footnotesize},
    title style={font=\small\bfseries, yshift=1pt},
    every axis plot/.append style={line width=1.1pt},
    clip=false,
    legend style={font=\footnotesize, draw=none, fill=none}
  },
  aqmcurve/.style={aqmaccent, smooth},
  aqmcurvebroken/.style={aqmbroken, smooth},
  aqmcurvemuted/.style={aqmmuted, smooth, line width=0.7pt, dashed},
  aqmguide/.style={aqmmuted, line width=0.6pt, dash pattern=on 2pt off 1.5pt}
}

% --- Figure float helper -----------------------------------------------------
% Every schematic file uses this so floats behave identically across the set.
% Usage:
%   \begin{aqmfigure}{fig:s01:ssb}{Caption text.}
%     <tikzpicture>
%   \end{aqmfigure}
\newenvironment{aqmfigure}[2]{%
  \begin{figure}[htbp]
  \centering
  \def\aqmfiglabel{#1}\def\aqmfigcaption{#2}%
  \begin{adjustbox}{max width=\linewidth}%
}{%
  \end{adjustbox}
  \caption{\small\aqmfigcaption}
  \label{\aqmfiglabel}
  \end{figure}
}
